\documentclass{amsart}
\usepackage{amssymb}
\usepackage{amsthm}
\usepackage{amscd}
\usepackage{graphics}
\newtheorem*{reftheorem}{Theorem}
\newtheorem*{reflemma}{Lemma}
\newtheorem{theorem}{Theorem}
\newtheorem{example}{Example}
\newtheorem{introexample}{Example}
\newtheorem{ML}{Main Lemma}
\newtheorem*{theorem 1}{Theorem 1}
\newtheorem*{theorem 2}{Theorem 2}
\newtheorem*{theorem 3}{Theorem 3}
\newtheorem{lemma}{Lemma}[section]
\newtheorem{corollary}{Corollary}[section]
\newtheorem{prop}{Proposition}[section]
\newtheorem{introtheorem}{Theorem}
\newtheorem{introlemma}{Lemma}
\newtheorem{introprop}{Proposition}
\newtheorem*{refcor}{Corollary}
\newtheorem*{refprop}{Proposition}
\newtheorem*{definition}{Definition}
\newtheorem*{conj}{Conjecture}
\newtheorem{remark}{Remark}
\newtheorem{question}{Question}
\newtheorem*{ques}{Question}
\newtheorem*{purlem}{Pursing Lemma}
\newtheorem{map}{Mapping}
\newtheorem*{mapv}{Variation of Mapping 3}
\numberwithin{equation}{section}
\raggedbottom

\newcommand{\Dil}{\textrm{Dil}}
\newcommand{\Vol}{\textrm{Vol}}
\newcommand{\RR}{\mathbb{R}}
\newcommand{\ZZ}{\mathbb{Z}}

\begin{document}

\title{A proposition for first proof: solution}

\maketitle

Problem statement.

Suppose $(M, g)$ and $(N,h)$ are piecewise smooth Riemannian manifolds.   If $f: M \rightarrow N$ is piecewise smooth,  we define the $k$-dilation of $f$ as the infimal $\Lambda$ so that for every $k$-dimensional submanifold $\Sigma \subset M$, 

$$ \Vol_k( f(\Sigma) ) \le \Lambda \Vol_k (\Sigma). $$

\noindent We write the $k$-dilation of $f$ as $\Dil_k(f)$.   Equivalently,  $\Dil_k(f) = \sup_{x \in M} \| \Lambda^k df \|$.

\begin{prop}  Prove that there is a constant $k > 0$ so that the following holds.

Suppose that $R \subset \RR^4$ is a 4-dimensional rectangle with side lengths $R_1 \le R_2 \le R_3 \le R_4$ and $S$ is a 4-dimensional rectangle with side lengths $S_1 \le S_2 \le S_3 \le S_4$.   Suppose that $f: (R, \partial R) \rightarrow (S, \partial S)$ is a piecewise smooth map with degree 1 and with $\Dil_2(f) \le 1$.   Suppose that $R_1 \le k S_1$.   Then either

$$ R_3 R_4 > k S_3 S_4,  or $$

$$ R_1 R_2 R_3 R_4 > k S_1 S_2^{1/2} S_3^{3/2} S_4. $$

\end{prop}




We write $A \lesssim B$ to mean that there is some positive constant $C$ so that $ A \le C B$.  We write $A \sim B$ to mean that $A \lesssim B$ and $B \lesssim A$.  

\subsection{Isoperimetric lemma}

The first main ingredient of our proof is an isoperimetric inequality for relative cycles in a rectangle.   Thie lemma appeared in \cite{G} as Theorem 3,  but since \cite{G} was not published we include the full proof here.

Let $R$ denote the n-dimensional rectangle $[0, R_1] \times ...
\times [0, R_n]$, where the dimensions are ordered so that $R_1
\le ... \le R_n$.    If $z$ is a relative integral k-cycle in $R$, the filling volume
of $z$ is the smallest volume of any relative (k+1)-chain $y$
with $\partial y = z$.  Let $I^k_R(V)$ denote the largest filling
volume of any k-dimensional relative integral cycle in $R$ with
volume at most $V$.

Remark: We use the following definition for volume.  If a chain
$C$ is given by $\sum c_i f_i$ where $c_i \in \mathbb{Z}$ and
$f_i$ is a Lipschitz map from the standard k-simplex to $S$, then
the volume of the chain $C$ is defined to be $\sum |c_i|
Vol(f_i^* Euc)$, where $Euc$ denotes the Euclidean metric on $R$. 
This quantity is also called the mass of $C$.  

\begin{lemma} \label{isoperimetriclemma0}
$I_R^0(V) \le R_1 V$
\end{lemma}

\begin{proof} When $k=0$, $z$ is just a weighted sum of
points $\sum c_i p(i)$, where $c_i \in \mathbb{Z}$ and $p(i)$ is
a point in $R$.  The volume of $z$ is defined to be $\sum |c_i|$. 
A point $p$ with coordinates $(p_1, ..., p_n)$ bounds a segment
$[0, p_1] \times \{ p_2 \} \times ... \times \{ p_n \}$. 
Applying this operation to each point $p(i)$ with multiplicity
$c_i$, we get a filling of $z$ with volume at most $R_1 Vol_0(z)$.
\end{proof}


\begin{lemma} \label{isoperimetriclemma} (Isoperimetric lemma) There are constants $c(n) > 0, C(n)$ so that the
following holds for any $1 \le k \le n-1$.

If $V \le c(n) R_1 ... R_k$, then write $V = c(n) R_1 ... R_j
\rho^{k-j}$ for some $0 \le j \le k-1$ and some $\rho$ in the
range $R_j \le \rho \le R_{j+1}$. (These conditions determine $j$
and $\rho$ uniquely.)

$$\textrm{Then } I^k_R(V) \le C(n) R_1 ... R_j \rho^{k-j+1}.$$

In any case, $I^k_R(V) \le C(n) R_{k+1} V$.

\end{lemma}

The proof of the isoperimetric lemma is based on using the Federer-Fleming deformation theorem at multiple scales.  We first recall a version of the Federer-Fleming deformation theorem.  


\begin{theorem} (Federer-Fleming) Suppose that $z$ is a relative k-cycle in $R$.
Consider a rectangular lattice inside $R$ with each side-length
roughly equal to $L$ (up to a factor of 2), for some $L \le R_1$,
and suppose that the boundary of $R$ lies in the (n-1)-skeleton
of the lattice.  Then there is another relative cycle $z'$ in $R$
contained in the k-skeleton of the lattice and obeying the 
following inequalities.

1. The volume of $z'$ is at most $C(n) |z|$.

2. The filling volume of $z' - z$ is at most $C(n) L |z|$.

\end{theorem}

See \cite{G3} for a proof of this form of the Federer-Fleming deformation theorem.  

\begin{proof} [Proof of Lemma \ref{isoperimetriclemma}]

We prove the isoperimetric inequality by
induction on $k$.    Lemma \ref{isoperimetriclemma0} serves as a base case.

If $k \ge 2$,  suppose by induction that the Lemma holds for $k-1$.  

Suppose that $z$ is a k-cycle
with volume $V$.  We proceed in two cases.  If $V \le c(n)
R_1^k$, then we select $L = C(n) V^{1/k} \le R_1$ and pick a
rectangular lattice with sidelengths roughly $L$ and with
$\partial R$ in the (n-1)-skeleton of the lattice. Then we use
Lemma 1.1 to move $z$ to a new relative cycle $z'$ with volume at
most $C(n) V$ lying in the k-skeleton of our lattice. Since $C(n)
V \le L^k$, the new cycle $z'$ is simply $0$.  Lemma 1.1 also
guarantees us a homology from $z$ to $z'$ with volume at most
$C(n) L V$, which is at most $C(n) V^{\frac{k+1}{k}}$.  This
upper bound is the one we needed to prove.

In the second case, we suppose that $V \ge c(n) R_1^k$.  In this
case, we select $L = R_1$ and pick a rectangular lattice with
sidelengths roughly $L$ and with $\partial R$ in the
(n-1)-skeleton of the lattice.  We pick the lattice so that each
lattice point has $x_1$ coordinate either $0$ or $R_1$. Then we
use Lemma 1.1 to move $z$ to a new relative cycle $z'$ with
volume at most $C(n) V$ lying in the k-skeleton of our lattice. 
The homology from $z$ to $z'$ has volume at most $C(n) R_1 V$. 
The cycle $z'$ need not be $0$, but it is a union of interior
k-faces of our lattice.  Each interior k-face has the form $[0,
R_1] \times ...$, and so the cycle $z'$ has the special form $z'
= [0, R_1] \times z_1$ for some relative cycle $z_1$ in the
(n-1)-dimensional rectangle $[0, R_2] \times ... \times [0,
R_n]$.  The cycle $z_1$ has volume at most $C(n) V / R_1$.

If $k = 1$,  then we find a filling of $z_1$ using Lemma \ref{isoperimetriclemma0},  and if $k > 1$,  then we find a filling
of $z_1$ by induction.  
Therefore, $z_1$ bounds a relative chain $C_1$ with a certain
volume bound that we calculate below.  Then $z'$ bounds $[0, R_1]
\times C_1$.  We will calculate that the volume of this filling
obeys the inequality stated in the theorem.

If $k=1$,  then we get the bound 

$$ \Vol_2([0,R_1] \times C_1) = R_1 \Vol_1(C_1) \le R_1 R_2 \Vol_0(z_1) \le C(n) R_2 V,$$

\noindent which is the desired bound in the case that $V \ge c(n) R_1$.

Now we turn to the case $k > 1$.  

If the volume of $z$ is at most $c(n) R_1 ... R_k$, then the
volume of $z_1$ is at most $c(n-1) R_2 ... R_k$.  If the volume
of $z$ is equal to $c(n) R_1 ... R_j \rho^{k-j}$ for some $\rho$ in
the range $R_j \le \rho \le R_{j+1}$, then the volume of $z_1$ is
roughly $R_2 ... R_j \rho^{(k-1) - (j-1)}$ for the same $\rho$. 
By induction, $z_1$ bounds a chain $C_1$ with volume at most
$C(n-1) R_2 ... R_j \rho^{k-j+1}$, and so $[0, R_1] \times C_1$
has volume at most $C(n-1) R_1 ... R_j \rho^{k-j+1}$.  Also, the
homology from $z$ to $z'$ has volume at most $C(n) R_1 (R_1 ...
R_j) \rho^{(k-j)} \le C(n) R_1 ... R_j \rho^{k-j+1}$. Therefore,
the filling volume of $z$ is at most $C(n) R_1 ... R_j
\rho^{k-j+1}$.  This gives the desired bound if the volume of $z$ is at most $c(n) R_1 ... R_k$.

Regardless of the volume of $z$,  $z_1$ bounds a k-chain $C_1$ of volume at most
$C(n-1) R_{k+1} V / R_1$. Hence $z' = [0,R_1] \times z_1$ bounds
a (k+1)-chain of volume at most $C(n-1) R_{k+1} V$.  The homology
from $z$ to $z'$ has volume at most $C(n) R_1 V$.  Therefore, the
filling volume of $z$ is at most $C(n) R_{k+1} V$. \end{proof}



\subsection{Linking invariants}


An important character in the proof will be a certain homotopy invariant that we call a linking invariant.   We now describe this invariant and review standard properties.

 Consider the pair $(S^2 \vee S^2, *)$,  where $*$ is the wedge point of $S^2 \vee S^2$.    Let $\Omega_1, \Omega_2$ be generators of $H^2(S^2 \vee S^2,  *)$ corresponding to the two copies of $S^2$,  and let $\omega_1, \omega_2$ be differential forms representing $\Omega_1,  \Omega_2$.   We choose $\omega_i$ so that it is supported in the $i^{th}$ $S^2$ outside of a neighborhood of $*$. 

Suppose that $(M,  \partial M)$ is a compact manifold with boundary,  and $h \in H_3(M,  \partial M; \ZZ)$.   Suppose that $g: (M,  \partial M) \rightarrow (S^2 \vee S^2,  *)$ is a piecewise smooth map with $g^*(\Omega_i) = 0$ in $H^2(M,  \partial M; \ZZ)$ for $i=1,2$.   By the de Rham theorem for cohomology with compact support,  it follows that there are piecewise smooth differential forms $\alpha_1,  \alpha_2$ on $(M, \partial M)$ with

\begin{itemize}

\item $\alpha_i$ is supported in interior of $M$.

\item $d \alpha_i = g^*(\omega_i)$

\end{itemize}

If $z$ is a relative 3-chain in $(M, \partial M)$ with $[z] = h$,  we define

\begin{equation} \label{defL}
L_h(g) := \int_z \alpha_1 \wedge g^*(\omega_2)
\end{equation}

If $(M,  \partial M)$ is an oriented 3-manifold with fundamental homology class $[M]$, we abbreviate $L(g) = L_{[M]}(g)$.  

\begin{lemma} \label{lemmalink} Assuming that $g^*(\Omega_i) = 0 \in H^2(M, \partial M)$ for $i = 1, 2$,  $L_h$ obeys the following properties.

\begin{enumerate}


\item $L_h(g)$ is independent of the choice of $z$

\item If $\beta_1$ is a 1-form on $(M, \partial M)$ with $d \beta_1 = g^*(\omega_1)$,  even if the support of $\beta_1$ includes the boundary of $M$,  we still have

$$ L_h(g) = \int_z \beta_1 \wedge g^*(\omega_2). $$

\item $L_h(g)$ is a homotopy invariant of $g$. 

\item If $\phi: (L, \partial L) \rightarrow (M, \partial M)$ and $\phi_*(h_L) = h_M$,  then

$$ L_{h_L}(g \circ \phi) = L_{h_M} (g).  $$

\end{enumerate}

\end{lemma}

\begin{proof}

 \begin{enumerate}

\item Note that $d ( \alpha_1 \wedge g^*(\omega_2)) = g^*(\omega_1) \wedge g^*(\omega_2) = g^*(\omega_1 \wedge \omega_2) = 0$.   If $[z_1] = [z_2] = h$,  then $z_1 - z_2 = \partial y$ for a relative 4-chain $y$,  and the equality follows from Stokes theorem.

\item If $d \beta_1 = d \alpha_1$,  then $\int_z \beta_1 \wedge g^*(\omega_2) - \int_z \alpha_1 \wedge g^*(\omega_2) = \int_z (\alpha_1 - \beta_1) \wedge d \alpha_2 =  \int_z d \left( (\alpha_1 - \beta_1) \wedge \alpha_2 \right)$.   Apply Stokes theorem and note that $(\alpha_1 - \beta_1) \wedge \alpha_2$ vanishes on $\partial z \subset \partial M$ because $\alpha_2$ vanishes there.

\item Suppose $G: (M \times [0,1],  \partial M \times [0,1]) \rightarrow (S^2 \vee S^2, *)$ is a homotopy between $g_0$ and $g_1$.    Let $Z = z \times [0,1]$.   Choose $A_i$ to be  forms on $(M \times [0,1],  \partial M \times [0,1])$,  supported in $Int(M) \times [0,1]$ with $d A_i = G^*(\omega_i)$.   We have $d( A_1 \wedge G^*(\omega_2)) = 0$ as in item 1,  and so by Stokes theorem on $Z$,  

$$L_h(g_0) =  \int_{z \times \{ 0 \} } A_1 \wedge G^*(\omega_2) = \int_{z \times \{ 1 \} } A_1 \wedge G^*(\omega_2) = L_h(g_1). $$

\item Choose $\alpha_{i, M}$ supported on $Int(M)$ so $d \alpha_{i,M} = g^*(\omega_i)$.   Then define $\alpha_{i,L} = \phi^*(\alpha_{i,M})$,  supported in $Int(L)$ with $d \alpha_{i,L} = (g \circ \phi)^{*}(\omega_2)$.   Choose $z_L$ with $[z_L] = h_L$, and set $z_M = \phi(z_L)$ so $[z_M] = h_M$.   Then 

$$L_{h_L}(g \circ \phi) = \int_{z_L} \alpha_{1,L} \wedge (g \circ \phi)^*(\omega_2) = \int_{z_M} \alpha_{1,M} \wedge g^*(\omega_2) = L_{h_M}(g).$$ 

\end{enumerate}


\end{proof}

This linking invariant is not always zero.    Let $Q_1$ be the unit 3-cube.   There is a piecewise smooth map $g_1: (Q_1,  \partial Q_1) \rightarrow (S^2 \vee S^2, *)$ with $L(g_1) = 1$.   This is a standard construction.   For details,  see Proposition 5.2 in \cite{GH}.

We equip $(S^2 \vee S^2,  *)$ with a metric where each $S^2$ is a round sphere of radius $S_1$.   Using this metric,  we can choose the differential forms $\omega_i$ so that 

\begin{equation} \| \omega_i \|_{L^\infty} \lesssim S_1^{-2}.  \end{equation}

Let $Q_{S_1}$ be a 3-cube of side length $S_1$.   By scaling $g_1$,  we can construct a map $g_{S_1}: (Q_{S_1},  \partial Q_{S_1}) \rightarrow (S^2 \vee S^2,  *)$ with $\Dil_1(g_{S_1}) \lesssim 1$ and $L(g_{S_1}) = 1$.  

The connection between $k$-dilation and linking invariants such as this one is explored in \cite{GH}. 

\subsection{The proof of the Proposition}

\begin{proof}

For each $(a,b) \in [0, S_3] \times [0, S_4]$,  let $\Sigma_{a, b} := [0, S_1] \times [0, S_2] \times \{a \} \times \{ b \} \subset S$.   By the transversality theorem,  for almost every $(a,b) \in  [0, S_3] \times [0, S_4]$,  the map $f$ is transverse to $\Sigma_{a,b}$ and so $f^{-1}(\Sigma_{a,b})$ is a piecewise smooth submanifold.   

Define a set $A$ of good values of $(a,b)$ by

\begin{equation} \label{abnice}  A := \{ (a,b) \in [(1/4) S_3, (3/4) S_3] \times [(1/4) S_4, (3/4) S_4]: f \textrm{ is transverse to } \Sigma_{a,b} \} . \end{equation}  

We will prove a lower bound for $\Vol_2(\Sigma_{a,b})$ for each $(a,b) \in A$.  Note that the measure of $A$ is $|A| \sim S_3 S_4$.  

Next we discuss the topology of the pair $(S,  \partial S \cup \Sigma_{a,b})$.   Note that $H_2(\partial S \cup \Sigma_{a,b}; \ZZ) = \ZZ$.   To build a generator,  note that $\partial \Sigma_{a,b}$ is homologically trivial in $\partial S$,  and let $y$ be a 2-chain in $\partial S$ with $\partial y = - \partial \Sigma_{a,b}$.   Then let $z_{2,  S} = y + \Sigma_{a,b}$,  which is a cycle in $\partial S \cup \Sigma_{a,b}$ representing a generator of $H_2(\partial S \cup \Sigma_{a,b})$,  which we call $h_{2,S}$.   

Also note that $H_3(S,  \partial S \cup \Sigma_{a,b}; \ZZ) = \ZZ$.   To build a generator,  let $z_{3,  S}$ be a relative 3-chain in $(S, \partial S)$ with $\partial z_{3,  S} = \Sigma_{a,b}$.    Then $z_{3,  S}$ is a relative 3-cycle in $(S,  \partial S \cup \Sigma_{a,b})$ representing a generator of $H_3(S,  \partial S \cup \Sigma_{a,b}; \ZZ)$,  which we call $h_{3, S}$.   Note that $\partial h_{3,S} = h_{2,S}$.

Next we build a topologically interesting map

$$ g: (S,  \partial S \cup \Sigma_{a,b} ) \rightarrow (S^2 \vee S^2, *). $$

Define $S' = [0,S_1] \times [0,  S_2] \times [0, S_3]$.    Let $r: [0, S_3] \times [0, S_4] \rightarrow [0, S_3]$ be given by 

$$ r(y_3, y_4) = \min(  10 ( |y_3 - a| + |y_4 - b| ),  S_3 ).  $$

Now let $g_1 (y_1, y_2, y_3, y_4) = (y_1, y_2,  r(y_3, y_4))$ so $g_1: S \rightarrow S'$.    Notice that $g_1(\Sigma_{a,b}) \subset [0,S_1] \times [0, S_2] \times \{0 \} \subset \partial S'$ and by \eqref{abnice},  $g_1(\partial S) \subset \partial S'$.   Therefore $g_1$ is a map of pairs $g_1: (S,  \partial S \cup \Sigma_{a,b}) \rightarrow (S',  \partial S')$.   Also $g_{1, *}(h_{3,S})$ is the fundamental homology class of $(S',  \partial S')$.   

Suppose that $Q_1,  ...,  Q_N \subset S'$ are disjoint cubes of side length $S_1$,  with $N \sim  \frac{S_2 S_3}{S_1^2}$.   For each cube $Q$,  recall that we defined above a map
$g_Q: (Q, \partial Q) \rightarrow (S^2 \vee S^2, *)$ with $L(g_Q) = 1$ and $\Dil_1(g_Q) \lesssim 1$.   Now we define $g_2$ so that for each $Q$,  $g_2 |_Q = g_Q$,  and so that $g_2(x) = *$ if $x$ is not in any $Q$.   Then $g_2: (S', \partial S') \rightarrow (S^2 \vee S^2, *)$ has $\Dil_1(g_2) \lesssim 1$.

Since $H^2(S',  \partial S') = 0$,  $g_2^*(\Omega_i) = 0$ for $i=1,2$ and $L(g_2)$ is well defined.   Next we compute $L(g_2)$.   Since $g_2 |_Q = g_Q$,  we see that $g_2^*(\omega_i) |_Q = g_Q^*(\omega_i)$.   For each $Q$,  choose $\alpha_{1,Q}$ with compact support in $Int(Q)$ so that $d \alpha_{1,Q} = g_Q^{*}(\omega_1)$.   Then define $\alpha_{1,2}$ so that $\alpha_{1,2} |_Q = \alpha_{1,Q}$ for each $Q$ and $\alpha_{1,2}$ vanishes outside of the union of the cubes $Q$.    Then

$$ L(g_2) = \int_{S'} \alpha_{1,2} \wedge g_2^*(\omega_2) = \sum_{Q} \int_Q \alpha_{1,Q} \wedge g_Q^*(\omega_2) = \sum_{Q} L(g_Q) = \sum_{Q \in S'} 1 = N \sim \frac{S_2 S_3}{S_1^2}. $$

Now we define $g = g_2 \circ g_1: (S,  \partial S \cup \Sigma_{a,b}) \rightarrow (S^2 \vee S^2,  *)$.    We have $\Dil_1(g) \lesssim 1$ and so $\Dil_2(g) \lesssim 1$.   Since $g_{1, *}(h_{3,S}) = [S']$,  we have

\begin{equation}  L_{h_{3,S}}(g)  = L(g_2) \sim \frac{S_2 S_3}{S_1^2}.   \end{equation}

For each $(a,b) \in A$, we view the map $f$ as a map of pairs $f: (R,  \partial R \cup f^{-1} (\Sigma_{a,b})) \rightarrow (S,  \partial S \cup \Sigma_{a,b})$.   Recall that for $(a,b) \in A$,  $f$ is transverse to $\Sigma_{a,b}$ and so $f^{-1}(\Sigma_{a,b})$ is a piecewise smooth submanifold.    

We can define a homology class $h_{2,R} \in H_2(R,  \partial R \cup f^{-1}(\Sigma_{a,b}); \ZZ)$ just as we defined $h_{2,S}$: find a 2-chain $y$ in $\partial R$ with $\partial y = - \partial f^{-1}(\Sigma_{a,b})$ and set $z_{2,R} = y + f^{-1}(\Sigma_{a,b})$.   Since $f: (R,  \partial R) \rightarrow (S,  \partial S)$ is degree 1,  it follows that $f: (f^{-1}(\Sigma_{a,b}),  \partial f^{-1}(\Sigma_{a,b})) \rightarrow (\Sigma_{a,b},  \partial \Sigma_{a,b})$ is also degree 1,  and so $f_*(h_{2,R}) = h_{2, S}$.  

Now suppose that $z_3$ is a relative 3-chain in $(R,  \partial R)$ with $\partial z_3 = f^{-1}(\Sigma_{a,b})$.   We can also view $z_3$ as a relative 3-cycle in $(R,  \partial R \cup f^{-1}(\Sigma_{a,b}))$.   We let $h_{3,R} = [z_3]$.   Then we have $\partial h_{3,R} = h_{2,R}$.   

If we view $f$ as a map of pairs,  $f: (R,  \partial R \cup f^{-1} (\Sigma_{a,b})) \rightarrow (S,  \partial S \cup \Sigma_{a,b})$,  then $f_*(h_{3,R}) = h_{3,S}$.   To see this,  note that  $\partial f_*(h_{3,R}) = f_*(h_{2,R}) = h_{2, S}.  $   But the $\partial$ map $H_3(S,  \partial S \cup \Sigma_{a,b}) \rightarrow H_2(\partial S \cup \Sigma_{a,b})$ is an isomorphism,  and the only class $h$ with $\partial h = h_{2,S}$ is $h = h_{3,S}$.   

Now we consider the map $F = g \circ f: (R, \partial R \cup f^{-1}(\Sigma_{a,b}))  \rightarrow (S^2 \vee S^2, *)$.   Since $\Dil_1(g) \lesssim 1$,  we have $\Dil_2(g) \lesssim 1$ and so $\Dil_2( F) \lesssim 1$.   By Lemma \ref{lemmalink},  we also have

\begin{equation} \label{linkF} L_{h_{3,R}}(F) = L_{h_{3,S}}(g) \sim \frac{ S_2 S_3}{S_1^2}.  \end{equation}

Next we give an upper bound for $L_{h_{3,R}}(F)$.   Recall that $z_3$ is any relative 3-chain in $(R,  \partial R)$ with $\partial z_3 = f^{-1}(\Sigma_{a,b})$,  which we can also view as a relative  3-cycle in $(R,   \partial R \cup f^{-1}(\Sigma_{a,b}))$ with $[z_3] = h_{3, R}$.   

Recall that $\| \omega_i \|_{L^\infty} \lesssim S_1^{-2}$.   Since $\Dil_2(F) \lesssim 1$,  we have

\begin{equation}
\| F^* (\omega_i) \| \lesssim S_1^{-2}
\end{equation}

Next we choose $\beta_1$ a 1-form on $R$ with $d \beta_1 = F^*(\omega_1)$.    We choose $\beta_1$ by integrating $F^*(\omega_1)$ in the $x_1$ direction.   Let $(x_1, x_2, x_3, x_4)$ be the standard coordinates on $R = [0,R_1] \times ... \times [0, R_4]$.   If we write 

$$F^*(\omega_1) = \sum_{j =2}^4 c_j(x) dx_1 \wedge dx_j + \sum_{2 \le j_1 < j_2 \le 4} c_{j_1, j_2}(x) d x_{j_1} \wedge d x_{j_2}, $$

then

$$ \beta_1 = \sum_{j=2}^4 \left( \int_{0}^{x_1} c_j(x_1',  x_2,  x_3, x_4) dx_1' \right) dx_j. $$

Since $F^*(\omega_1)$ is closed,  it is a standard fact that $d \beta_1 = F^*(\omega_1)$ -- see the proof of the Poincare lemma in Chapter 1 of \cite{BT}.    We have $\| c_j \|_\infty \le \|F^*(\omega_i) \|_{L^\infty} \lesssim S_1^{-2}$,  and $0 \le x_1 \le R_1$,  and so

\begin{equation} \label{betabound}
\| \beta_1 \|_{\infty} \lesssim R_1 S_1^{-2}
\end{equation}

Now we have

$$ |L_{h_{3,R}}(F)| = \left| \int_{z_3} \beta_1 \wedge F^*(\omega_2) \right| \lesssim \Vol_3(z_3) R_1 S_1^{-4}.  $$

Combining this with \eqref{linkF} and rearranging, we get

\begin{equation} \label{stepSz}  R_1^{-1} S_1^2 S_2 S_3 \lesssim \Vol_3(z_3).  \end{equation}

Next we upper bound $\Vol_3(z_3)$ using the isoperimetric lemma,  Lemma \ref{isoperimetriclemma}.   For reference,  we record here the relevant special cases of Lemma \ref{isoperimetriclemma}.   

\begin{lemma} \label{lemmaisoperimetric2} There are constants $c > 0, C$ so that the
following holds.

If $y$ is a relative 2-cycle in $R$ with 

\begin{equation} \label{casea} c R_1^2 \le \Vol_2(y) \le c R_1 R_2,  \end{equation}

then $y = \partial z$ for a relative 3-chain $z$ in $R$ with

\begin{equation} \label{bounda}  \Vol_3(z) \le C R_1^{-1} \Vol_2(y)^2.  \end{equation}

If $y$ is a relative 2-cycle in $R$ with

\begin{equation} \label{caseb} \Vol_2(y) \le c R_1^2,  \end{equation}

then $y = \partial z$ for a relative 3-chain $z$ in $R$ with

\begin{equation} \label{boundb}  \Vol_3(z) \le C \Vol_2(y)^{3/2} \le C R_1^3.  \end{equation}

\end{lemma}

We apply this theorem with $y = f^{-1}(\Sigma_{a,b})$ and $z = z_3$.   Since $R_1 \le k S_1$,  we have $\Vol_3(z_3) \gtrsim R_1^{-1} S_1^2 S_2 S_3 \ge k^{-3} R_1^3$.   By choosing $k > 0$ small enough,  we can guarantee that \eqref{boundb} does not hold.   Examining Lemma \ref{lemmaisoperimetric2},  we see that only two possibilities remian:for every $(a,b) \in A$,  either

\begin{equation} \label{casei} \Vol_2(f^{-1}(\Sigma_{a,b})) \gtrsim R_1 R_2,  or \end{equation}

\begin{equation} \label{caseii} \Vol_2[f^{-1}(\Sigma_{a,b})] \gtrsim R_1^{1/2} \Vol_3(z_3)^{1/2} \gtrsim S_1 S_2^{1/2} S_3^{1/2} \end{equation}

Now using $\Dil_2(f) \le 1$ and the coarea inequality,  we have

$$ R_1 R_2 R_3 R_4 = \Vol_4 (R) \ge \int_{A} \Vol_2(f^{-1}(\Sigma_{a,b})) da db. $$

Using $|A| \sim S_3 S_4$ and plugging in the bound from \eqref{casei} and \eqref{caseii},  we see that

$$ R_1 R_2 R_3 R_4 \gtrsim R_1 R_2 S_3 S_4,  or $$

$$ R_1 R_2 R_3 R_4 \gtrsim S_1 S_2^{1/2} S_3^{3/2} S_4. $$


\end{proof}




\begin{thebibliography}{99}

\bibitem{BT} R. Bott and L. Tu,  {\it Differential Forms in Algebraic Topology}

\bibitem{G3} Guth, L., Notes on Gromov's systolic estimate,
Geometriae Dedicata 123, Number 1, Dec. 2006, 113-129.

\bibitem{G} Guth, L., Area-expanding embeddings of rectangles, arXiv:0710.0403 

\bibitem{GH} Guth, L.  Isoperimetric inequalities and rational homotopy invariants, arXiv:0802.3550 


\end{thebibliography}


\end{document}


